\chapter{트랜스포머 블록 조립하기}
\label{ch:transformer_block}

이 장에서는 5장의 어텐션과 4장의 임베딩을 결합하여 \keyterm{트랜스포머 블록(transformer block)}을 만든다. 어텐션만으로는 부족하고, FFN(Feed-Forward Network), 잔차 연결(residual connection), 레이어 정규화(layer normalization)가 함께 있어야 비로소 의미 있는 표현을 학습할 수 있다.

\section{트랜스포머 블록의 구조}

하나의 트랜스포머 블록은 다음 네 가지 컴포넌트로 이루어진다:

\begin{enumerate}
  \item \textbf{Multi-Head Attention}: 토큰 간 정보 교환
  \item \textbf{Feed-Forward Network (FFN)}: 위치별 비선형 변환
  \item \textbf{Residual Connection}: 입력을 출력에 더하여 깊은 네트워크 학습
  \item \textbf{Layer Normalization}: 학습 안정화
\end{enumerate}

이 네 가지가 어떻게 조합되는지가 트랜스포머 블록의 핵심이다. 1권에서는 Pre-LN 구조(GPT-2 스타일)를 사용한다.

\begin{figure}[htbp]
\centering
\begin{tikzpicture}[
  every node/.style={font=\small\sffamily},
  box/.style={draw=inkblue, fill=inkblue!8, rounded corners=2pt, minimum width=3cm, minimum height=0.7cm},
  norm/.style={draw=inkgreen, fill=inkgreen!10, rounded corners=2pt, minimum width=2.5cm, minimum height=0.6cm},
  sum/.style={draw=inkred, fill=inkred!10, circle, minimum size=0.6cm}
]
  \node[box] (in) at (0, 0) {input $x$};
  \node[norm] (ln1) at (0, -1) {LayerNorm};
  \node[box] (attn) at (0, -2) {Multi-Head Attention};
  \node[sum] (add1) at (0, -3) {$+$};
  \node[norm] (ln2) at (0, -4) {LayerNorm};
  \node[box] (ffn) at (0, -5) {FFN (GELU)};
  \node[sum] (add2) at (0, -6) {$+$};
  \node[box] (out) at (0, -7) {output};

  % residual 우회선
  \draw[inkblue, rounded corners=4pt] (-2, 0) -- (-2, -3);
  \draw[-{Stealth[length=4pt]}, inkblue] (-2, -3) -- (add1);
  \draw[inkblue, rounded corners=4pt] (2, -3) -- (2, -6);
  \draw[-{Stealth[length=4pt]}, inkblue] (2, -6) -- (add2);

  % 메인 흐름
  \draw[-{Stealth[length=4pt]}, inkblue] (in) -- (ln1);
  \draw[-{Stealth[length=4pt]}, inkblue] (ln1) -- (attn);
  \draw[-{Stealth[length=4pt]}, inkblue] (attn) -- (add1);
  \draw[-{Stealth[length=4pt]}, inkblue] (add1) -- (ln2);
  \draw[-{Stealth[length=4pt]}, inkblue] (ln2) -- (ffn);
  \draw[-{Stealth[length=4pt]}, inkblue] (ffn) -- (add2);
  \draw[-{Stealth[length=4pt]}, inkblue] (add2) -- (out);

  % 라벨
  \node[inkgray, font=\scriptsize, anchor=west] at (-3.5, -1.5) {residual};
  \node[inkgray, font=\scriptsize, anchor=west] at (2.5, -4.5) {residual};
\end{tikzpicture}
\caption{Pre-LN 트랜스포머 블록 구조. 어텐션과 FFN 각각에 잔차 연결과 LayerNorm이 있음.}
\label{fig:transformer-block}
\end{figure}

\section{Feed-Forward Network (FFN)}

어텐션은 토큰 간 ``수평적'' 정보 교환을 수행한다. 반면 FFN은 각 토큰 위치에서 ``수직적'' 비선형 변환을 수행한다. 구조는 단순하다:

\[
\text{FFN}(x) = \text{Linear}_2(\text{GELU}(\text{Linear}_1(x)))
\]

$\text{Linear}_1$은 $d_{model} \to d_{ff}$로 확장하고, $\text{Linear}_2$는 $d_{ff} \to d_{model}$로 다시 축소한다. $d_{ff}$는 보통 $4 d_{model}$이며, 이 ``확장-축소'' 패턴이 비선형 표현력을 준다.

\subsection{직관: 왜 확장-축소인가?}

$d_{model} = 512$를 $d_{ff} = 2048$로 확장했다가 다시 512로 축소하는 이유는 ``더 큰 표현 공간에서 연산''하기 위해서다. 고차원 공간에서는 더 복잡한 결정 경계를 그릴 수 있다. 활성화 함수(GELU)의 비선형성과 결합하여, FFN은 단순한 선형 변환 이상의 것을 학습한다.

실제로 FFN은 ``키-값 메모리'' 역할을 한다는 연구 결과(Geva et al. 2021)도 있다. $\text{Linear}_1$의 행벡터가 ``키'', $\text{Linear}_2$의 열벡터가 ``값''으로 동작하여, FFN이 사실 정보 검색 메커니즘이라는 해석이다.

\begin{lstlisting}[language=Python]
import torch.nn.functional as F

class FeedForward(nn.Module):
    def __init__(self, d_model: int, d_ff: int, dropout: float = 0.1):
        super().__init__()
        self.w1 = nn.Linear(d_model, d_ff)
        self.w2 = nn.Linear(d_ff, d_model)
        self.dropout = nn.Dropout(dropout)
        self._init_weights()

    def _init_weights(self) -> None:
        nn.init.normal_(self.w1.weight, mean=0.0, std=0.02)
        nn.init.normal_(self.w2.weight, mean=0.0, std=0.02)
        if self.w1.bias is not None:
            nn.init.zeros_(self.w1.bias)
        if self.w2.bias is not None:
            nn.init.zeros_(self.w2.bias)

    def forward(self, x: torch.Tensor) -> torch.Tensor:
        """x: [batch, seq, d_model] -> [batch, seq, d_model]"""
        return self.dropout(self.w2(F.gelu(self.w1(x))))
\end{lstlisting}

\section{GELU vs ReLU}

원본 트랜스포머는 ReLU를 사용했지만, GPT-2 이후로는 \keyterm{GELU(Gaussian Error Linear Unit)}가 표준이다. GELU는 ReLU의 ``냉혹한'' 0/1 임계점 대신 가우시안 누적 분포를 사용하여 부드럽게 전환한다.

\begin{formula}[GELU]
\[
\text{GELU}(x) = x \cdot \Phi(x) \approx 0.5 x \left(1 + \tanh\left[\sqrt{2/\pi}(x + 0.044715 x^3)\right]\right)
\]
$\Phi(x)$는 표준정규분포의 누적분포함수. $x < 0$에서도 작은 음수값을 허용하여 그래디언트 흐름이 더 좋다.
\end{formula}

\begin{figure}[htbp]
\centering
\begin{tikzpicture}[
  scale=0.8,
  every node/.style={font=\small\sffamily}
]
  % 축
  \draw[inkblue, line width=0.5pt, ->] (-3, 0) -- (3, 0) node[right] {$x$};
  \draw[inkblue, line width=0.5pt, ->] (0, -0.5) -- (0, 2.5) node[above] {activation};
  % ReLU (꺾임)
  \draw[inkred, line width=1pt] (-3, 0) -- (0, 0) -- (2.5, 2.3);
  \node[inkred, font=\scriptsize] at (2, 2.3) {ReLU};
  % GELU (부드러운 곡선)
  \draw[inkblue, line width=1pt, domain=-3:2.5, samples=50] plot (\x, {0.5*\x*(1+tanh(0.7978*(\x+0.044715*\x*\x*\x)))});
  \node[inkblue, font=\scriptsize] at (1.5, 1.5) {GELU};
\end{tikzpicture}
\caption{ReLU vs GELU. ReLU는 $x < 0$에서 정확히 0이지만, GELU는 작은 음수값을 허용하여 부드럽게 전환.}
\label{fig:gelu-relu}
\end{figure}

\subsection{왜 GELU가 더 나은가?}

ReLU는 $x < 0$에서 완전히 0이 되어 그래디언트도 0이 된다 (``죽은 ReLU'' 문제). GELU는 작은 음수값을 허용하여 그래디언트가 항상 흐른다. 또한 전환점이 부드러워 미분 가능하여 최적화가 안정적이다.

실제로 GPT-2, BERT, RoBERTa 등 대부분의 현대 모델이 GELU를 사용한다. 2권에서는 SwiGLU라는 더 발전된 활성화 함수를 다룬다.

\section{잔차 연결 (Residual Connection)}

\keyterm{잔차 연결(residual connection)}은 입력을 출력에 직접 더하는 연산이다: $y = x + f(x)$. 이 단순한 아이디어가 깊은 네트워크 학습을 가능하게 했다. He et al. (2016)의 ResNet이 이미지 분류에서 100층 이상 학습을 가능하게 한 것이 그 시작이다.

\subsection{직관: 왜 잔차가 도움인가?}

잔차가 없으면 네트워크는 $f(x) = y$를 학습해야 한다. 깊어질수록 이 매핑을 학습하기 어렵다. 잔차가 있으면 $f(x) = y - x$, 즉 ``잔차''만 학습하면 된다. 입력 $x$가 이미 $y$와 가깝다면 $f$는 작은 보정만 하면 된다.

또한 역전파 시 그래디언트가 ``지름길''을 통해 직접 흐른다. $y = x + f(x)$의 미분은 $\frac{dy}{dx} = 1 + f'(x)$다. $f'(x) = 0$이어도 그래디언트는 1이 되어 사라지지 않는다.

트랜스포머 블록에는 두 개의 잔차 연결이 있다: 어텐션 주변 하나, FFN 주변 하나. GPT는 12$\sim$96층이므로 잔차 없이는 학습이 불가능하다.

\begin{warnbox}[잔차 없는 100층 네트워크]
실험적으로 잔차를 빼고 100층 네트워크를 학습하면 손실이 줄어들지 않는다. 그래디언트가 역전파 도중 사라져 앞쪽 층이 학습되지 않기 때문이다. 잔차 연결은 단순하지만 깊은 네트워크의 필수 요소다.
\end{warnbox}

\section{Pre-LN vs Post-LN}

LayerNorm의 위치에 따라 두 가지 변형이 있다.

\subsection{Post-LN (Vaswani 원본)}

원본 트랜스포머는 LayerNorm을 어텐션/FFN \emph{뒤}에 배치했다:

\[
x = \text{LN}(x + \text{Attn}(x))
\]

하지만 이 구조는 깊은 네트워크에서 학습이 불안정하다. 초기화에 매우 민감하고 warmup이 없으면 발산하기 쉽다.

\subsection{Pre-LN (GPT-2 스타일)}

LayerNorm을 어텐션/FFN \emph{앞}에 배치한다:

\[
x = x + \text{Attn}(\text{LN}(x))
\]

이 구조는 잔차 경로에 LayerNorm이 없어 그래디언트가 더 잘 흐른다. 깊은 네트워크에서도 warmup 없이 안정적으로 학습되어 현대 LLM의 표준이다.

\subsection{왜 Pre-LN이 더 안정적인가?}

Post-LN에서 잔차 경로 $x \to x + \text{Attn}(x) \to \text{LN}(\cdot)$를 거치면, 그래디언트가 LayerNorm을 통과해야 한다. LayerNorm은 분산으로 나누는 연산이라 그래디언트를 스케일링하는데, 깊은 네트워크에서 이 스케일링이 누적되면 그래디언트가 사라질 수 있다.

Pre-LN은 잔차 경로 $x \to x + \text{Attn}(\text{LN}(x))$에서 $x$가 LayerNorm을 거치지 않고 직접 흐른다. 어텐션 출력만 LayerNorm을 거친 후 더해진다. 그래디언트가 identity 경로를 통해 그대로 흘러 깊은 네트워크에서도 학습이 가능하다.

\begin{lstlisting}[language=Python]
class TransformerBlock(nn.Module):
    def __init__(self, config: ModelConfig, layer_norm_style: str = "pre"):
        super().__init__()
        self.layer_norm_style = layer_norm_style
        self.ln1 = nn.LayerNorm(config.d_model, eps=config.layer_norm_eps)
        self.ln2 = nn.LayerNorm(config.d_model, eps=config.layer_norm_eps)
        self.attn = MultiHeadAttention(config)
        self.ffn = FeedForward(config.d_model, config.d_ff, config.dropout)

    def forward(self, x, mask=None):
        if self.layer_norm_style == "pre":
            # Pre-LN: 정규화 -> 어텐션 -> 잔차
            x = x + self.attn(self.ln1(x), mask=mask)
            x = x + self.ffn(self.ln2(x))
        else:
            # Post-LN: 어텐션 -> 잔차 -> 정규화
            x = self.ln1(x + self.attn(x, mask=mask))
            x = self.ln2(x + self.ffn(x))
        return x
\end{lstlisting}

\begin{warnbox}[Pre-LN을 써야 하는 이유]
Post-LN은 초기 그래디언트가 잔차 경로가 아닌 LayerNorm을 통해서만 흐르기 때문에, 깊은 네트워크에서 사라질 수 있다. Pre-LN은 잔차 경로에 아무것도 없어 그래디언트가 identity하게 흘러 100층 이상에서도 학습이 가능하다. 단, Pre-LN은 마지막에 추가 LayerNorm이 필요하다 (다음 장에서 \code{ln\_f}로 구현).
\end{warnbox}

\section{잔차 스케일링}

GPT-2는 잔차 스트림에 직접 기여하는 가중치(어텐션의 $W_O$, FFN의 $W_2$)를 더 작게 초기화한다. 층이 깊어질수록 잔차 출력의 분산이 커지는 것을 방지하기 위함이다.

\begin{lstlisting}[language=Python]
def _init_residuals(self, module: nn.Module) -> None:
    """잔차 스트림에 직접 기여하는 프로젝션 가중치를 작게 초기화."""
    if isinstance(module, nn.Linear):
        if module.weight.shape[0] == self.config.d_model:
            # 잔차로 들어가는 출력: 더 작은 표준편차로 초기화
            nn.init.normal_(module.weight, mean=0.0,
                            std=0.02 / math.sqrt(2 * self.config.n_layers))
\end{lstlisting}

$1/\sqrt{2N}$ 스케일링은 $N$층일 때 잔차 출력의 분산이 층 수에 비례해 커지는 것을 보정한다. $N$이 클수록 더 작은 초기화가 필요하다.

\subsection{왜 $2N$인가?}

트랜스포머 블록마다 잔차에 기여하는 출력이 2개 (어텐션, FFN)다. $N$개 블록이면 총 $2N$개의 출력이 잔차에 더해진다. 각 출력의 분산이 $\sigma^2$이면 합의 분산은 $2N \sigma^2$이다. 이것이 1이 되려면 $\sigma = 1/\sqrt{2N}$이 되어야 한다.

\section{블록 테스트}

\begin{lstlisting}[language=Python]
def test_block_shape(small_config):
    block = TransformerBlock(small_config)
    x = torch.randn(2, 8, small_config.d_model)
    out = block(x)
    assert out.shape == (2, 8, small_config.d_model)

def test_block_residual_connection(small_config):
    """잔차 연결이 있으면 입력과 출력의 차이가 유한해야 함."""
    block = TransformerBlock(small_config)
    block.eval()
    x = torch.randn(2, 8, small_config.d_model)
    out = block(x)
    diff = (out - x).abs().mean()
    assert diff.item() < 5.0  # 잔차가 너무 크면 정규화 문제

def test_feedforward_shape(small_config):
    ffn = FeedForward(small_config.d_model, small_config.d_ff)
    x = torch.randn(2, 8, small_config.d_model)
    out = ffn(x)
    assert out.shape == x.shape
\end{lstlisting}

\section{트랜스포머 블록의 계산량}

블록 하나의 FLOPs는 대략:
\begin{itemize}
  \item QKV 프로젝션: $3 \cdot L \cdot d^2$
  \item 어텐션 점수: $L^2 \cdot d$
  \item 어텐션 출력 프로젝션: $L \cdot d^2$
  \item FFN: $2 \cdot L \cdot d \cdot d_{ff} = 8 L d^2$ (since $d_{ff} = 4d$)
\end{itemize}

총 $O(L \cdot d^2)$이며, 어텐션의 $O(L^2 d)$는 $L < d$일 때 무시할 수 있다. 실제 GPT-3 ($d = 12288$, $L = 2048$)에서 $L \cdot d^2 = 3 \times 10^{11}$, $L^2 \cdot d = 5 \times 10^{10}$이라 FFN이 더 지배적이다.

\section{요약}

\begin{itemize}
  \item 트랜스포머 블록 = 어텐션 + FFN + 잔차 연결 + LayerNorm.
  \item FFN은 위치별 비선형 변환 (Linear $\to$ GELU $\to$ Linear).
  \item GELU는 ReLU보다 부드러운 전환으로 학습 안정성 향상.
  \item 잔차 연결은 깊은 네트워크 학습의 핵심.
  \item Pre-LN이 Post-LN보다 학습 안정성이 좋아 현대 LLM 표준.
  \item 잔차 스케일링 $1/\sqrt{2N}$로 깊은 네트워크의 분산 폭발을 방지.
\end{itemize}

\section{연습 문제}

\begin{enumerate}
  \item FFN의 $d_{ff}$를 $4d_{model}$ 대신 $d_{model}$로 하면 어떤 일이 일어나는가?
  \item 잔차 연결을 제거하고 12층 모델을 학습하면 어떤 현상이 관찰되는가?
  \item Pre-LN에서 마지막 블록 후에 추가 LayerNorm이 필요한 이유는?
  \item GELU 대신 Sigmoid Linear Unit (SiLU, $x \cdot \sigma(x)$)를 사용하면 어떤 차이가 있는가?
  \item 잔차 스케일링 없이 100층 모델을 초기화하면 학습 초기에 어떤 일이 일어나는가?
\end{enumerate}

다음 장에서는 이 블록을 $N$층 쌓아 올려 완전한 GPT 모델을 만들고, 학습 루프를 구축한다.
